% Capítulo 6: Conclusiones

\chapter{Conclusiones}

\label{Chapter6}

\lhead{Capítulo 6. \emph{Conclusiones}}

\textit{[Recupera los objetivos iniciales y explica si se han cumplido y por qué. Debes comenzar con el primer objetivo específico y terminar con el general, explicando cómo y por qué has cumplido con el 100\% de lo que te planteaste al inicio del proyecto.]}

\textit{[Se recomienda un tamaño de 5 a 10 páginas para este capítulo.]}

%----------------------------------------------------------------------------------------

\section{Limitaciones}

\textit{[Errores cometidos y limitaciones de los resultados obtenidos, es decir, qué ha fallado para que tu proyecto no fuera perfecto.]}

%----------------------------------------------------------------------------------------

\section{Prospectiva}

\textit{[¿Cómo vas a continuar o completar este proyecto? Aunque nunca lo hagas, debes decir cómo deberías hacerlo.]}

%----------------------------------------------------------------------------------------

\section{Consideraciones finales}

\textit{[Reflexión personal sobre las competencias adquiridas, lo que se ha aprendido a lo largo de todo el Máster y con la elaboración del TFM, tanto a nivel de contenidos como personal. Autoevaluación y propuestas de mejora. Trata de responder a las siguientes cuestiones:]}

\begin{itemize}
    \item \textit{[Reflexión final sobre los resultados obtenidos en el proyecto a todos los niveles.]}
    \item \textit{[Reflexión personal sobre las competencias adquiridas a lo largo del TFM, tanto a nivel profesional como del método científico.]}
    \item \textit{[Reflexión personal sobre lo que se ha aprendido a lo largo de todo el Máster, tanto a nivel profesional como personal.]}
    \item \textit{[Autoevaluación y propuestas de mejora.]}
    \item \textit{[Agradecimientos a todos los que han hecho posible el proyecto.]}
    \item \textit{[Deseo positivo final relacionado con el ámbito al que se ha dedicado el trabajo como colofón a todo el proyecto. Ejemplo: ``No puedo terminar este proyecto sin reflejar mi deseo de que la situación de X mejore en el futuro, para lo cual será muy importante Y. Por mi parte espero que esta pequeña contribución que ha supuesto este TFM sea el inicio de Z\ldots'']}
\end{itemize}
